\subsection{Weighted canonical-flat root fibres with triple points}

In this subsection the marked weighted data include the reduced
central analytic space, its component multiplicities, the central
log space of the normalized cyclic root, and the old cone cells
before their semistable subdivision.  A coarsening is an actual
proper toric morphism between the displayed monoid charts.  These
data retain more than the dual incidence complex.

\begin{statement}{C.2.UN.76}{Proved}{Exact weighted toric descent and compact occurrence}
Let \(f:X\to\Delta_t\) be a proper degeneration of connected complex
two-tori, with smooth total space and full fibre
\(F=\sum_i m_iV_i\) whose reduced support is simple normal crossing.
Suppose there is a triple point and that the actual orders of a
generator of \(f_*K_X\) satisfy
\[
 \operatorname{div}(\Omega)=\sum_i k_iV_i,\qquad
 \frac{k_i+1}{m_i}=\lambda\quad\hbox{for every }i.
\tag{C.WF.1}
\]
Then the following conclusions hold.

The normalized cyclic root, followed by the explicit height-one
subdivision below, is a positive canonical-flat toric degeneration
of UN71.  The original marked weighted support is exactly its
retained coarse-cone quotient by the root group.  Thus this branch
is described by the toric data \((N,\Gamma,\mathscr P,c)\) of UN71,
the retained periodic coarse triangles, and the actual cyclic
monomial maps (C.WF.9)--(C.WF.11), with their specified smooth
coarse invariant charts (C.WF.28).

There are precisely two possible common ratios:
\[
 \lambda=1,\quad k_i=m_i-1;
 \qquad\hbox{or}\qquad
 \lambda=\tfrac12,\quad k_i=\tfrac{m_i}{2}-1.
\tag{C.WF.2}
\]
In the second case every \(m_i\) is even.  Every positive-dimensional
open stratum of the original support has Euler number zero.
Consequently
\[
             e(F_{\mathrm{red}})
               =\#\{\hbox{original triple points}\}>0.
\tag{C.WF.3}
\]
In particular this entire weighted branch contains no zero-Euler
support.

Every such marked weighted log support has a realization in a
smooth compact complex threefold over \(\PP^1\), with connected
two-torus fibres away from two values.  The two marked weighted
supports are the given one and its complex conjugate.  The actual
old-cone contractions and finite invariant-ring maps are retained
at both ends, so the reduced analytic support, all its gluing maps,
and every full-fibre multiplicity are unchanged.  The total Euler
number of this compact realization is twice (C.WF.3).

The exact original nonreduced fibre and its entire neighbourhood
are recovered by the original coarsening and quotient morphisms.
The compact occurrence assertion preserves the marked reduced
log support and multiplicity cycle; it does not identify an
arbitrary original higher infinitesimal thickening or full
degeneration germ with the new compact realization.
\end{statement}

\begin{proof}
\par\medskip\noindent\textbf{The actual normalized-root and coarse-cone maps.}
Choose \(M\) divisible by all \(m_i\), and use \(t=s^M\).
Locally a holomorphic root of the nonzero unit in \(f\) gives
actual coordinates in which
\[
 t=x_1^{m_1}\cdots x_r^{m_r},\qquad 1\le r\le3.
\tag{C.WF.4}
\]
The coordinate change is invertible because its derivative in
the changed normal direction is a unit.  Keep its choice and
all its transition units in the normalization diagram.
The normalized root has \(\gcd(m_1,\ldots,m_r)\) local components.
Each is the normal toric chart with cocharacter lattice and cone
\[
 L=\{(v,h)\in\Z^r\oplus\Z:\ \textstyle\sum_i m_iv_i=Mh\},
 \qquad \sigma=\{v_i\ge0\}.
\tag{C.WF.5}
\]
The remaining \(3-r\) coordinates are unchanged.  Its old rays
are \((M e_i/m_i,1)\), which are primitive because their height
is one.  Formula (C.WF.5), its saturated dual monoid, and the
characters \(x_i,s\) are the actual normalization maps; no
unramified branches are discarded.

The height-one slice is an interval or a lattice triangle.
Subdivide its edges at every lattice point and triangulate
using all its lattice points.  Every resulting empty lattice
triangle is unimodular: its Euclidean lattice area is \(1/2\)
by the elementary lattice-area formula.  Equivalently its
three lifted height-one vectors are an integral basis of \(L\).
Thus the resulting charts are literally
\(\C^3\), with \(s=u_1u_2u_3\), or their two-factor and
one-factor versions.

For the actual original charts, perform this refinement by
successive height-one star subdivisions.  Each is the normalized
blowup of its monomial ideal.  Changing a boundary equation by a
holomorphic unit multiplies each monomial generator by a unit
and leaves that ideal unchanged.  The normalized blowup maps
therefore glue with the original tangent-dependent transition
units retained.  We have not treated those original overlaps
as constant toric maps.  After all height-one lattice points
have been inserted the same empty-triangle argument gives the
unimodular refinement.

The subdivisions can be chosen equivariantly.  A stabilizer of an old
coarse cone in the actual cyclic root fixes its old rays
individually: they lie over the distinct original SNC
components at that stratum.  It therefore acts trivially on
the height-one simplex.  Choose one subdivision on each orbit
of old triangles after the common edge subdivisions, and
translate it to the other members of the orbit.  Equivalently
perform the corresponding star operations on their whole
orbits; the old component labels and the trivial simplex
stabilizers give the same restrictions on every shared face.
The proper normalized blowup maps are compatible on overlaps.  They give
an actual proper equivariant morphism
\[
 \widetilde X\longrightarrow X_M
       :=\operatorname{Norm}(X\times_{\Delta_t}\Delta_s),
 \qquad X_M/\mu_M=X.
\tag{C.WF.6}
\]
The equality on the right is the actual finite analytic
quotient: invariant meromorphic functions are those on \(X\),
and normality of the smooth \(X\) makes its integral invariant
algebra exactly \(\OO_X\).  The retained coarsening in
(C.WF.6) is essential; taking the quotient of the subdivided
space alone would usually retain additional divisors.
The semistable support is embedded SNC globally.  Its old
rays retain the distinct original component labels.  A new
ray retains its original stratum and its lattice position
in that stratum's old simplex.  Transition maps and root
inertia fix these positions on each old cone.  Rays occurring
together in a subdivided cone therefore cannot become two
branches of one global prime.  The star charts of each
prime glue as the smooth toric boundary charts of that
labelled ray.

\par\medskip\noindent\textbf{The canonical orders, with their exact normalization.}
At a height-one ray \((v,1)\), direct differentiation of
(C.WF.4) gives
\[
 \operatorname{ord}_{v}(\rho^*\Omega)
       =\sum_i(k_i+1)v_i-1
       =M\lambda-1.
\tag{C.WF.7}
\]
For an old ray this is \(M(k_i+1)/m_i-1\).
For every new ray it is the affine interpolation of those
old values.  This also proves the formula in a two-factor
chart after wedging the unchanged tangent differential.
The pulled-back nonzero coefficient is a unit throughout
each normalization chart.  Hence
\[
       \widetilde\Omega
          =s^{-(M\lambda-1)}\rho^*\Omega
\tag{C.WF.8}
\]
is a nowhere-zero holomorphic three-form on \(\widetilde X\).
Its exponent is an integer by (C.WF.7).
An old triple cone contains at least one unimodular triangle,
so the root fibre has a triple point.  UN68 and UN71 apply to
this actual reduced semistable root.  In particular its marked
central log space is the toric space of
\((N,\Gamma,\mathscr P,c)\).  The old cones give a periodic
coarsening \(\mathscr P_0\) of this same developed triangulation.

Notice the normalization of (C.WF.1).  In a weighted fibre it
is not necessary that some \(k_i\) be zero.  Positivity of
the \(k_i+1\) gives \(\lambda>0\).  If \(\lambda>1\), every
integer \(k_i-m_i=(\lambda-1)m_i-1\) is nonnegative, so
\(\Omega/t\) is holomorphic on the smooth \(X\).  This
contradicts that \(\Omega\) generates \(f_*K_X\).
Thus \(0<\lambda\le1\).

\par\medskip\noindent\textbf{The exact central finite action.}
Write \(\zeta=\exp(2\pi i/M)\).  Lift the root generator
to the developed affine plane and to the path-product
trivialization of the toric cocycle in UN71.  The resulting
maps on the open torus are
\[
 g(z,s)=(a\,s^{b_0}z^A,\zeta s),\qquad
 A\in\operatorname{GL}(N),\quad b_0\in N,\quad a\in T_N,
 \qquad p\longmapsto Ap+b_0.
\tag{C.WF.9}
\]
They preserve both \(\mathscr P\) and the marked coarsening
\(\mathscr P_0\).  Here \(z^A\) denotes the map whose
pullback on characters is \(A^{\mathsf T}\).

We justify the central comparison, including its constants.
Each component map preserves its full toric boundary, so on
its dense torus it is a torus translation followed by the
integral map of its fan.  This follows directly by pulling
back a character: its divisor is a specified integral
combination of boundary primes, and the quotient of two
functions with that divisor is a holomorphic unit on the
compact component, hence a constant.  In the global
characters at a vertex \(p\), its dense component coordinates
are \(z^m s^{-\langle m,p\rangle}\).  Substitution in
(C.WF.9) gives the component constant
\(a\,\zeta^{-(Ap+b_0)}\).
Conversely the actual double-curve map fixes the tangential
character comparison, and its marked normal-product log
section fixes the other character comparison.  In the
three-factor chart these are exactly the two substitutions
just displayed.  Thus after the UN71 path-product
trivialization the quantities
\(a_p\zeta^{Ap+b_0}\) agree across every edge and equal one
constant \(a\).  This proves (C.WF.9) on the entire central
log atlas.  It is a central-log comparison; no conjugacy of
the arbitrary original whole germ has been asserted.

The lifted \(g\) need not have order \(M\) on this developed
cover.  Instead it normalizes the actual period maps
\(g_\gamma(z,s)=(c_\gamma s^\gamma z,s)\) and satisfies
\[
 \begin{split}
 A\Gamma&=\Gamma,\qquad A^M=I,\\
 c_\gamma^A&=\zeta^{A\gamma}c_{A\gamma},\\
 g^M&=g_{\gamma_0}\quad\hbox{for a retained }\gamma_0\in\Gamma .
 \end{split}
\tag{C.WF.10}
\]
All are equalities of monomial maps.  For clarity the last
one means
\[
 \gamma_0=\sum_{j=0}^{M-1}A^{M-1-j}b_0,\qquad
 \prod_{j=0}^{M-1}(a^{A^{M-1-j}})
       \ \zeta^{\sum_{j=0}^{M-1}jA^{M-1-j}b_0}
       =c_{\gamma_0}.
\tag{C.WF.11}
\]
The central log comparison preserves these relations.
Indeed two monomial maps with the same affine map and base
character differ by a constant torus translation; their
actions on a height-one component determine that constant.
Equality of the actual central maps therefore implies the
displayed equality on every toric chart.

The relative residue of (C.WF.8) transforms under the root
generator by \(\zeta^{-M\lambda}\).  On the developed toric
log atlas it is a common nonzero scalar times
\(d\log z_1\wedge d\log z_2\), and (C.WF.9) transforms this
residue by \(\det A\).  The scalar is common because the
developed transition maps are translations and the double
residues agree.  Consequently
\[
                  \zeta^{-M\lambda}=\det A\in\{1,-1\}.
\tag{C.WF.12}
\]
Together with \(0<\lambda\le1\), this proves (C.WF.2),
including the evenness of every \(m_i\) in the second case.

\par\medskip\noindent\textbf{Euler number on the actual original strata.}
Let \(V_i^\circ\) be an original component with all double
curves removed.  The strict root components above it have
finite unramified maps to \(V_i^\circ\); locally these are
exactly the branches of \(x_i^{m_i}=s^M\) with the
nonzero units retained.  Their total degree is \(m_i\).
The semistable subdivisions change nothing there.
On every old root component this inverse image is its full
toric complement \((\C^*)^2\).  The sum of their Euler
numbers is therefore zero, and finite-cover multiplicativity
gives \(e_c(V_i^\circ)=0\).

For a connected original double curve \(C\), remove its
original triple points and write \(C^\circ\).  The normalized
root over this open curve is locally the two-dimensional
monomial root times \(C^\circ\).  Its subdivision is the
corresponding chain times that curve.  It creates no triple
point over the interior of \(C^\circ\).  Choose a boundary
curve of an old strict root component which dominates \(C\);
boundary curves supported over an original triple point
are not used here.  Every such dominating curve is a toric
\(\PP^1\), and the
preimage of the removed original endpoints is precisely
its two toric endpoints.  Thus its open part is \(\C^*\)
and maps finitely and unramifiedly to \(C^\circ\).
This remains true component by component when the root
cover is disconnected over \(C^\circ\).  Its positive total
degree gives \(e_c(C^\circ)=0\).  Additivity on the original
SNC stratification now proves (C.WF.3).  This argument
uses the retained old cells, rather than the Euler number
of the subdivided root fibre.

\par\medskip\noindent\textbf{The period equations retaining the finite action.}
Choose \(N=\Z^2\), \(\Gamma=Q\Z^2\) and logarithms \(C\)
of \(c\), as in UN73.  Initially put
\[
 \Lambda= \Z^2,\quad \Lambda_0=Q^{-1}\Z^2,\quad
 B=(2\pi i)^{-1}Q^{-1}C,\quad
 D=Q^{-1}AQ,\quad
 d=Q^{-1}b_0,\quad
 a_0=(2\pi i)^{-1}Q^{-1}\log a .
\tag{C.WF.13}
\]
The additive action is \(w\mapsto Dw+d\,l+a_0\),
where \(l=\log s/(2\pi i)\).  Replace these real lattices
by their actual images under one real invertible matrix
which makes the finite-order \(D\) orthogonal for an
averaged positive metric.  Replace \(B,D\) by their
conjugates and \(d,a_0\) by their images under the same
matrix.  We retain the names
\(\Lambda\subset\Lambda_0\), \(B,D,d,a_0\) for these
actual images.  In particular no integral lattice is
replaced by the abstract standard lattice in this step.
Let \(T\) be the actual real matrix taking this vector
coordinate back to \((2\pi i)^{-1}\log z\).

For \(\ell\in\Lambda\), (C.WF.10) becomes
\[
       (DB-BD-D/M)\ell\in\Lambda_0.
\tag{C.WF.14}
\]
Define
\[
 d_M=\sum_{j=0}^{M-1}D^{M-1-j}d\in\Lambda,\quad
 a_M=\sum_{j=0}^{M-1}D^{M-1-j}a_0,\quad
 e_M=\sum_{j=0}^{M-1}jD^{M-1-j}d .
\]
The order relation is precisely
\[
                   a_M+e_M/M-Bd_M\in\Lambda_0 .
\tag{C.WF.15}
\]
These are lattice-valued identities obtained by taking
the chosen actual logarithms of (C.WF.10)--(C.WF.11).

\par\medskip\noindent\textbf{An equivariant compact middle family.}
In the orthogonal coordinates put
\[
 J=\begin{pmatrix}0&1\\-1&0\end{pmatrix},\qquad
 DJ=(\det D)JD.
\]
Choose \(k\ge1\) with \(\zeta^k=\det D\).
Choose \(R>0\) so that the symmetric part \(S\) of
\(\operatorname{Im}B+RI\) is positive definite, and put
\[
 B'=B+iRI,\quad a'_0=a_0+iRd,\quad
 \kappa=e^{-2\pi R},\qquad
 V_s(q)=q+s^kJ\overline q,\quad H_s=I+s^kJ.
\tag{C.WF.16}
\]
The simultaneous shift of \(a_0\) is required: it is
exactly the substitution \(u=\kappa s\) in the original
affine map.  Equations (C.WF.14)--(C.WF.15) remain true
with \(B',a'_0\), since their added terms cancel.

On the whole punctured line use the actual period lattice
\[
 \mathcal L_s=V_s(\Lambda_0)+
       \{\,V_s(B'\ell)+l(s)V_s(\ell):\ell\in\Lambda\,\}.
\tag{C.WF.17}
\]
Changing the logarithm adds an angular period.  The
determinant calculation of UN73 applies in the chosen
real coordinates, with a fixed positive basis factor.
Writing \(\operatorname{Im}B'=S+aJ\),
\(r=|s|^k\), \(x=\operatorname{Re}(s^k)\) and
\(y=-\log|s|/(2\pi)\), its variable factor is
\[
 (1+r^2)^2\det S+(1-r^2)(1+r^2)y\operatorname{tr}S
 +y^2(1-r^2)^2+\bigl((1+r^2)a+2xy\bigr)^2>0 .
\tag{C.WF.18}
\]
Thus (C.WF.17) gives a proper smooth torus family over
\(\C^*\), including the points where \(H_s\) is singular.

The finite action on this family is the actual map
\[
 (v,s)\longmapsto
 \bigl(Dv+V_{\zeta s}(a'_0)
                  +l(s)V_{\zeta s}(d),\ \zeta s\bigr).
\tag{C.WF.19}
\]
The identity \(DV_s=V_{\zeta s}D\) follows from the
choice of \(k\).  Equation (C.WF.14) proves that
(C.WF.19) preserves the actual period subgroups.
Iterating it \(M\) times gives the period
\[
 V_s(a'_M+e_M/M)+l(s)V_s(d_M)
       =V_s(B'd_M+n)+l(s)V_s(d_M),\qquad n\in\Lambda_0,
\tag{C.WF.20}
\]
by (C.WF.15).  Hence its induced action has order
dividing \(M\).  Its base action has order \(M\), so
the action is faithful and free over \(\C^*\).
All these computations use only real coefficients
\(j/M\) inside \(V_s\); the complex number \(l(s)\)
is kept outside.  No complex-linearity of \(V_s\)
is presumed.

\par\medskip\noindent\textbf{The exact zero-end filling and original action.}
Near zero use
\[
                    z=\exp(2\pi iT H_s^{-1}v).
\tag{C.WF.21}
\]
This divides by exactly the angular lattice.
The other periods give the variable toric holonomies
of UN73, with central values \(\kappa^\gamma c_\gamma\).
In this coordinate (C.WF.19) is
\[
 (z,s)\longmapsto
   \bigl(a(s)s^{b_0}z^A,\zeta s\bigr),\qquad
 a(s)=\exp\bigl(2\pi iT H_{\zeta s}^{-1}
                              V_{\zeta s}(a'_0)\bigr).
\tag{C.WF.22}
\]
The coefficient is a holomorphic nonzero vector and
\(a(0)=\kappa^{b_0}a\).
Using the actual local base \(u=\kappa s\)
therefore restores exactly \(a\), \(c_\gamma\),
and the given central root action.  Both
\(\mathscr P\) and its marked old coarsening are
preserved.  Apply the variable-unit toric chart
construction of UN73, first with the subdivision
and then with its proper coarsening.  The monomial
maps (C.WF.22) extend to both atlases and commute
with the actual coarsening morphism.

\par\medskip\noindent\textbf{The reflected and conjugate infinity end.}
Set \(q=s^{-1}\), replace the vector frame by
\(v_\infty=q^kv\), and put
\[
 H^\infty_q=q^kI+J,\qquad
 V^\infty_q(w)=q^kw+J\overline w .
\]
Reverse the two nonangular period generators as
in UN73.  The new period group has generators
\[
 H^\infty_q\Lambda_0,\qquad
 -V^\infty_q(B'\ell)+l(q)H^\infty_q\ell
                 \quad(\ell\in\Lambda).
\tag{C.WF.23}
\]
The exact finite map in this frame is
\[
 (v_\infty,q)\longmapsto
 \bigl((\det D)Dv_\infty+
       V^\infty_{\zeta^{-1}q}(a'_0)
       -l(q)V^\infty_{\zeta^{-1}q}(d),\
       \zeta^{-1}q\bigr).
\tag{C.WF.24}
\]
Since \((\det D)D H^\infty_q=
H^\infty_{\zeta^{-1}q}D\), the exponential
coordinate with frame \(H^\infty_q\) has central
affine vertex map \(p\mapsto Ap-b_0\).
Use the reflected old and subdivided fans
\(-\mathscr P_0\) and \(-\mathscr P\).
After the local base change \(v_{\mathrm{base}}=\kappa q\),
the central map is exactly
\[
 (z_\infty,v_{\mathrm{base}})
  \longmapsto
  (\overline a^{-1}v_{\mathrm{base}}^{-b_0}
                   z_\infty^A,\ \zeta^{-1}v_{\mathrm{base}}),
\tag{C.WF.25}
\]
and its period constants are \(\overline c_\gamma\).
Here \(\overline{a'_0}=\overline{a_0}-iRd\)
and \(-d\,l(q)=-d\,l(v_{\mathrm{base}})+iRd\),
so the recorded shifts cancel exactly.

Finally apply the actual torus inversion
\(z_\infty\mapsto z_\infty^{-1}\) and relabel the
period group by \(\gamma\mapsto-\gamma\).
This sends the reflected fans to the original
ones, sends the map (C.WF.25) to the conjugate
of (C.WF.9), and leaves the conjugate period
constants correct because
\(\overline c_{-\gamma}^{-1}=\overline c_\gamma\).
Thus the full coarse and subdivided central
log spaces, including the finite action, are
the complex conjugates of the prescribed data.
The two base coordinates satisfy
\(u\,v_{\mathrm{base}}=\kappa^2\).

\par\medskip\noindent\textbf{Finite stabilizers and smooth coarse quotients.}
We check the local smoothness assertion rather than
inferring it from an abstract quotient.  At a point
of a retained coarse central chart, take its finite
point stabilizer after quotienting by \(\Gamma\).
The actual Kummer chart (C.WF.4)--(C.WF.5) shows that
its inertia fixes each normal ray and is the identity
on the tangent stratum: the latter cover is unramified.
This property is retained by the exact central log
comparison.  Its monomial lattice map is therefore
identity on the ray span and on its quotient, and,
being of finite order, is the identity on the whole
lattice.  We use a sufficiently small invariant
neighbourhood of this point.  A stabilizer of a
lower-dimensional cone which moves points on its
torus stratum is not used as a point stabilizer.
In the newly constructed chart the point stabilizer's
action on a character \(\chi^m\)
is multiplication by a holomorphic unit \(c_m(s)\).
These coefficients depend only on \(s\), by
(C.WF.22) or (C.WF.24), and are multiplicative in \(m\).
Their central constants are the original stabilizer
characters.  For a generator of order \(d\), with
base character \(\xi\), let
\[
 L_m(s)=\log\!\frac{c_m(s)}{c_m(0)},\qquad
 B_m(s)=\frac1d\sum_{j=0}^{d-1}jL_m(\xi^js).
\tag{C.WF.26}
\]
The logarithm is the unique one vanishing at zero.
The finite-order identity gives
\(\sum_jL_m(\xi^js)=0\), and hence
\[
 B_m(\xi s)-B_m(s)=L_m(s).
\tag{C.WF.27}
\]
The coordinate change
\(\chi^m\mapsto e^{-B_m(s)}\chi^m\)
therefore makes every stabilizer coefficient constant.
The functions \(B_m\) are additive in \(m\), so this is
an actual automorphism of the saturated monoid algebra,
with the displayed inverse.  The base character has
constant coefficient, so its \(B_m\) is zero and the
change fixes \(s\).  Averaging has only been used for
this finite chart stabilizer, not for the developed
lift with \(g^M=g_{\gamma_0}\).

We can now use precisely the original coarse invariant
monoid chart, retained in (C.WF.5)--(C.WF.6):
\[
 \left(\OO_{X_M,\mathrm{coarse}}\right)^{\mathrm{Stab}}
        =\C\{x_1,x_2,x_3\},\qquad
       s^M=x_1^{m_1}\cdots x_r^{m_r}.
\tag{C.WF.28}
\]
The same equation holds in the new chart after the
explicit unit change (C.WF.26).  This proves smoothness
of the quotient total space and its exact component
multiplicities.  It also proves the embedded SNC
condition for its reduced fibre, including the original
component labels.  Reduction and finite invariants
commute in characteristic zero, so the central
coarsening and quotient identify the actual original
reduced analytic support and all its gluing maps.

\par\medskip\noindent\textbf{Compactness and the precise retained scope.}
The determinant (C.WF.18) and real fundamental
parallelepipeds give a proper middle torus family
on every compact annulus.  The bounded-unit
toric filling argument of UN73 gives proper
fillings at both ends.  A toric coarsening has
the same fan support and is proper; the coarse
fillings therefore remain proper.  The exact
exponential maps above identify them with the
middle family on punctured discs.  They glue
to a Hausdorff compact normal root space over
\(\PP^1_s\).  Its finite action is the actual
(C.WF.19), extended by the displayed monomial
maps at both ends.  Its quotient is proper
over \(\PP^1_s/\mu_M=\PP^1\) and is smooth
by (C.WF.28); it is free over the punctured
base, where each quotient fibre is one of
the actual connected two-tori.

The two marked reduced supports and their
multiplicities are the prescribed one and
its conjugate.  Euler additivity and the
zero Euler number of a smooth torus give
the stated total Euler number.  On the
original family, (C.WF.6) recovers its full
nonreduced fibre by its actual invariant
analytic algebra, not just by a cycle.
The constructed compact family instead
uses the explicit variable period units
in (C.WF.17).  Their higher \(s\)-jets were
not identified with those of an arbitrary
original germ.  No such identification is
needed for the stated compact occurrence
of the marked reduced log support and
full multiplicity cycle.  This completes
all the asserted comparisons.
\end{proof}
